\thispagestyle{empty}
\begin{center}
\includegraphics[width=1.0in]{cbit.png}

\vspace*{0.3in}
\fontsize{12pt}{12pt}\selectfont \textbf{DEPARTMENT OF ARTIFICIAL INTELLIGENCE AND DATA SCIENCE}\\[0.2cm]
\fontsize{14pt}{14pt}\selectfont\textbf{CHAITANYA BHARATHI INSTITUTE OF TECHNOLOGY}\\[0.2cm]
\textbf{HYDERABAD -- 500075}\\
\vspace*{0.3in}
\fontsize{16pt}{16pt}\textbf{\underline{PROJECT -- WK MAPPING}}
\end{center}

% Rows as written by the project team (MP_main_report.docx).
\fontsize{10pt}{13pt}\selectfont
\begin{table}[H]
    \centering
    \renewcommand{\arraystretch}{1.3}
    \begin{tabular}{|c|p{11cm}|}
        \hline
        \textbf{WK} & \textbf{How the project addresses it} \\
        \hline
        WK2 & Statistics and data analysis: Cohen's $\kappa$, 95\% Wilson intervals, precision, recall, F1, coverage and conditional accuracy, and confusion matrices on the 45 development cases; TF-IDF lexical similarity for concept clustering. \\
        \hline
        WK3 & Computing fundamentals: a fixed eight-label decision schema, deterministic evidence rules, SHA-256 hashing of frozen inputs and gates, and a replay cache keyed by prompt, schema, model and input hash. \\
        \hline
        WK4 & Specialist knowledge at the forefront of the discipline: a locally served language model (Qwen2.5 7B Instruct) prompted for typed relations, ESCO/SKOS taxonomy relations, and abstention through a rule-based reject option. \\
        \hline
        WK6 & Engineering practice with current tools: Python, Ollama, scikit-learn, pdfplumber, pandas, FastAPI and PostgreSQL, Flutter, a blinded Google Form, and a tested back end (pytest). \\
        \hline
        WK7 & Responsible use of AI where results inform curriculum decisions: no model decision is applied until a pre-declared gate is met on independent human labels, so the model stays in advisory-only Safe Mode. \\
        \hline
        WK8 & Engagement with current research on semantic similarity, ontology matching, ESCO skills research, inter-rater agreement, language-model annotation and selective classification, used to evaluate the module critically. \\
        \hline
        WK9 & Honest reporting: label sources are kept apart, AI-assisted answers are not presented as independent human validation, the reviewer form is blinded, and the gate is never relaxed after results are known. \\
        \hline
    \end{tabular}
    \caption*{Mapping of the Project with Knowledge and Attitude Profile}
\end{table}

\newpage
\thispagestyle{empty}
\begin{center}
\includegraphics[width=1.0in]{cbit.png}

\vspace*{0.3in}
\fontsize{12pt}{12pt}\selectfont \textbf{DEPARTMENT OF ARTIFICIAL INTELLIGENCE AND DATA SCIENCE}\\[0.2cm]
\fontsize{14pt}{14pt}\selectfont\textbf{CHAITANYA BHARATHI INSTITUTE OF TECHNOLOGY}\\[0.2cm]
\textbf{HYDERABAD -- 500075}\\
\vspace*{0.3in}
\fontsize{16pt}{16pt}\selectfont\textbf{\underline{PROJECT -- SDG MAPPING}}
\end{center}

\fontsize{10pt}{13pt}\selectfont
\begin{table}[H]
    \centering
    \renewcommand{\arraystretch}{1.3}
    \begin{tabular}{|c|p{11cm}|}
        \hline
        \textbf{SDG} & \textbf{Relevance of the project} \\
        \hline
        SDG4 & The platform helps universities see how far a curriculum covers the skills employers ask for and how it changes between versions. The module evaluated here makes the term-to-concept judgements behind those coverage results checkable, so that curriculum review can rest on evidence rather than on unverified model output. \\
        \hline
        SDG8 & Matching curriculum topics to the skills named in job postings, through the shared ESCO vocabulary, supports the alignment of what graduates learn with labour-market demand. The contribution is indirect: this report evaluates the reliability of that matching, not employment outcomes. \\
        \hline
    \end{tabular}
    \caption*{Mapping of the Project with Sustainable Development Goals}
\end{table}

\newpage
